
\documentclass[a4paper,twoside,openany]{article}
\usepackage[utf8]{inputenc}
\usepackage[T1]{fontenc}
\usepackage[french]{babel}
\usepackage{geometry}
\usepackage{graphicx} % Required for inserting images



\usepackage{marginnote}
%\usepackage{ulem} %pour les caractères barrés

\usepackage{multirow}
\usepackage{xifthen}

\usepackage{makeidx}
\usepackage{imakeidx}

\geometry{hmargin=3cm,vmargin=2cm}
\linespread{1.15}
\title{Transcription de l'Affranchissement Général du mandement de Beaufort en Savoie}
\author{Nicolas Morane - étudiant M1 SDH Panthéon-Sorbonne}
\date{15 Janvier 2025}

\usepackage[backend=biber, citestyle=verbose-trad2, bibstyle=verbose-trad2, citepages=omit, ibidpage=true, idemtracker=constrict, hyperref=false, sorting=nyt]{biblatex}
\bibliography{bibliographie.bib}




\makeindex[name=nom, title=Index nominum, options= -s perso.ist, columns=1]

\usepackage{glossaries}
\makeglossaries
\newglossaryentry{rière}
{
    name=rière,
    description={"dans le territoire de"}
}
\newglossaryentry{cens}
{
    name=cens,
    description={Le cens est dû par celui qui a la propriété utile d'un bien foncier (généralement un paysan) à celui qui en a la propriété éminente (le seigneur, qui à l'Époque moderne n'est pas nécessairement un noble)}
}
\newglossaryentry{servis}
{
    name=servis,
    description={"Les servis, appelés aussi censes, correspondent à des redevances annuelles et doivent être
payés perpétuellement au seigneur, en sus des frais d’acquisition (introge) d’une terre
dépendante de la seigneurie. Ils peuvent être dûs en argent, mais souvent en nature, redevances
en blé par exemple, ou autres denrées. Les servis pèsent parfois très lourds eu égard à la
rentabilité de la terre. Ces redevances seigneuriales, s’apparentant à une sorte de bail
emphytéotique, s’ajoutent aux tailles, gabelles, dîmes et autres impôts"\footnote{\cite{Gachet2022}}}
}
\newglossaryentry{laodsindemnité}
{
    name=laods d'indemnité,
    description={Laods d'indemnité est ...}
}
\newglossaryentry{laods}
{
    name=laods,
    description={"Les loads sont des droits de mutations perçus par un seigneur lors de la vente d’une terre par
un de ses sujets. Ils sont payés par l’acheteur. S’ils constituent une protection contre la
spéculation foncière des plus riches, ils sont aussi un frein aux acquisitions par les plus pauvres,
aux remembrements et à l’optimisation du travail des terres. Ils peuvent se monter jusqu’à un
sixième de la valeur des transactions.
Les loads et les servis concernent toutes les personnes qui ont acquis des biens fonciers
dépendant du fief d’un seigneur. A noter que le seigneur peut être une personne physique mais
également une abbaye, un évêché etc. Etre dépendant d’un fief ne signifie pas dans tous les cas
une dépendance personnelle envers le seigneur, ni forcément l’appartenance à une classe sociale
inférieure. Ceux qui dépendent personnellement d’un seigneur et qui sont de catégorie sociale
inférieure sont les personnes dites taillables, mainmortables ou corvéables à merci. On utilise
peu le terme de serf à l’époque moderne en Savoie. Sur ces personnes pèsent des charges
seigneuriales supplémentaires, notamment des corvées, et surtout le droit d’échute." \footnote{\cite{Gachet2022}} }
}

\newglossaryentry{suffertes}
{
    name=suffertes,
    description={droit de mutation que l'on percevait en sus du laod quand un fonds taillable était vendu à un homme de condition non taillable ou quand un fonds franc était acheté par un homme taillable}
}
\newglossaryentry{dîmes}
{
    name=dîmes,
    description={redevance en nature ou en argent, portant principalement sur les revenus agricoles collectés en faveur de l’Église}
}
\newglossaryentry{décimes}
{
    name=décimes,
    description={voir dîme}
}
\newglossaryentry{corvées}
{
    name=corvées,
    description={"Les corvées sont des prestations périodiques dues aux seigneurs par leurs sujets. Il peut s’agir
de labours de terres, de coupes de bois, de transports de matériaux etc. Selon leur nature elles
peuvent être intégrées aux servis sans faire l’objet d’une terminologie à part" \footnote{\cite{Gachet2022}}}
}
\newglossaryentry{alpéages}
{
    name=alpéages,
    description={appelé aussi "Auciège" est une redevance en nature c'était  le fruit d'un ou plusieurs  jours de lait pour l'utilisation d'un alpage en montagne. pour une étude sur l'Auciège voir \cite{Duparc}}
}
\newglossaryentry{droits de haut siège}
{
    name=droits de haut siège,
    description={à definir}
}
\newglossaryentry{favettiers}
{
    name=favettiers,
    description={(Savoie) (Droit) Paysans, bourgeois, gentilhommes assujettis aux redevances seigneuriales (load, sufferte, plaid, corvée, cens). }
}
\newglossaryentry{égances}
{
    name=égances,
    description={fractionnement des charges foncières lors du partage d’une censive}
}
\newglossaryentry{péages}
{
    name=péages,
    description={péages est ...}
}


\newglossaryentry{bans champêtres}
{
    name=bans champêtres,
    description={bans de justice, voir aussi \cite{Berthier}}
}

\newglossaryentry{remèdes possessoriaux}
{
    name=remèdes possessoriaux,
    description={remèdes possessoriaux est ...}
}

\newglossaryentry{bénéfices d'instances}
{
    name=bénéfices d'instances,
    description={bénéfices d'instances est ...}
}

\newglossaryentry{rescindants}
{
    name=rescindants,
    description={Demande tendante à faire annuler un acte, un jugement. On a jugé le rescindant. Par cet arrêt, on n’a jugé que le rescindant. Nous avons gagné le rescindant, c’est une présomption en notre faveur pour le rescisoire}
}

\newglossaryentry{rescisoire}
{
    name=rescisoire,
    description={L’objet principal pour lequel on s’est pourvu, soit contre un acte, soit contre un jugement, et qui reste à juger, quand l’acte ou le jugement a été annulé. Le rescindant et le rescisoire ne sont pas jugés par le même arrêt}
}
\newglossaryentry{droit de pignoratitie}
{
    name=droit de pignoratitie,
    description={Pignoratif ? Qui est relatif à un gage. Endossement pignoratif, qui transfère à l’endossataire les droits d’un créancier gagiste. Vieilli. Contrat pignoratif, forme de contrat de vente à réméré qui, dissimulant un prêt sur gage et s’apparentant à un pacte commissoire, est prohibée par la loi.}
}
\newglossaryentry{nomination des officiers de justice}
{
    name=nomination des officiers de justice,
    description={le chatelain (souvent un notaire) est choisi par le seigneur pour le représenter sur son fief qui lui délègue notamment le rôle d'officier de basse justice}
}

\newglossaryentry{greffe}
{
    name=greffe,
    description={greffe est ...}
}
\newglossaryentry{leyde sur le bétail}
{
    name=leyde sur le bétail,
    description={La leyde était sous l'Ancien Régime en France (surtout au Moyen Âge) un droit féodal, sous la forme d'impôt, levé sur les marchandises, denrées et bestiaux vendus en foire et marché}
}


\newglossaryentry{hale}
{
    name=hale,
    description={à définir}
}


\newglossaryentry{mense archiépiscopale}
{
    name=mense archiépiscopale,
    description={ revenus liés aux possessions du clergé. On parle de Mense curiale, monacale, abbatiale, épiscopale. C’est un revenu destiné à la subsistance des ecclésiastiques, aux besoins du culte, à l’entretien des biens et de l’assistance aux pauvres.}
}
\newglossaryentry{droit des langues dû par les bouchers}
{
    name=droit des langues dû par les bouchers,
    description={est un "prélèvement en nature des langues de boeufs et de vaches, et même parfois de tous les autres animaux abattus par les bouchers, ce droit est le plus souvent remplacé par une indemnité annuelle compensatrice" cf G Sangnier}
}
\newglossaryentry{recherche des cristaux}
{
    name=recherche des cristaux,
    description={droit coutumier local, le seigneur se réservant les cristaux trouvés sur son territoire}
}
\newglossaryentry{droit d’avoir un banc placé}
{
    name=droit d’avoir un banc placé,
    description={droit d’avoir un banc placé est ...}
}
\newglossaryentry{placer des patibulaires avec ses armoiries}
{
    name=placer des patibulaires avec ses armoiries,
    description={est un moyen pour le seigneur d'un lieu de marquer son pouvoir de justice (haute, moyenne ou basse}
}
\newglossaryentry{échuttes}
{
    name=échuttes,
    description={"Dans les procès et les débats qui ont régulièrement eu lieu entre les promoteurs
d’affranchissements des « mainmortables » et leurs détracteurs, c’est sur le droit d’échute
semble-t-il qu’on s’oppose le plus radicalement. En quoi consiste-t-il ? C’est le droit de
propriété qu’un seigneur peut avoir sur les biens de ses sujets « taillables », lorsque ceux-ci
décèdent sans héritiers « directs » (enfants mâles en général). Des coutumes régionales peuvent
différer : en certains lieux ce ne sera pas la totalité des biens mais une partie seulement, en
d’autres endroits l’existence de filles au décès pourra aussi moduler le droit d’échute, etc.
Comme souvent sous l’ancien régime, pas de règle universelle sans particularités locales. Mais
l’esprit est là : impossible à un serf de disposer et de tester librement de ses biens. D’où l’origine
possible du mot « mainmortable », qui aurait la main morte, incapable de coucher par écrit un
testament.
Ce droit d’échute est d’autant plus mal vécu qu’il s’applique même sur des biens acquis à
l’extérieur de la seigneurie. Ces mainmortables – taillables ou serfs – font également l’objet de
mesures vexatoires : certaines fonctions ou certaines tenues vestimentaires leurs sont interdites." \footnote{\cite{Gachet2022}} }
}

\newglossaryentry{terriers}
{
    name=terriers,
    description={registre contenant les lois et usages d'une seigneurie, la description des bien-fonds, les droits et conditions des personnes, ainsi que les redevances et obligations auxquelles elles sont soumises}
}
\newglossaryentry{litterés}
{
    name=litterés,
    description={autres écrits}
}
\newglossaryentry{doages}
{
    name=doages,
    description={doages est ...}
}
\newglossaryentry{banc du droit}
{
    name=banc du droit,
    description={banc du droit est ...}
}
\newglossaryentry{prône}
{
    name=prône,
    description={Instruction, accompagnée d'avis, qu'un prêtre fait aux fidèles à la messe paroissiale du dimanche}
}
\newglossaryentry{droit de reachapt}
{
    name=droit de reachapt,
    description={droit de réachat}
}
\newglossaryentry{Laide}
{
    name=Laide,
    description={ pour Leyde ? La leyde était sous l'Ancien Régime en France (surtout au Moyen Âge) un droit féodal, sous la forme d'impôt, levé sur les marchandises, denrées et bestiaux vendus en foire et marché}
}
\newglossaryentry{droit de pêche de chasse}
{
    name=droit de pêche de chasse,
    description={droit de pêche, de chassee est ...}
}

















\begin{document}

\maketitle

\section{Introduction}

Le document transcrit ci-dessous est conservé au sein des archives départementales de Savoie (73), dans les archives communales de la commune d'Hauteluce sous la cote : "E.suppl.944 - DD 1 (Cahier) - in 4° 20 feuillets"

Cela correspond à un document de 40 pages dont 3 vierges, la transcription proprement dite est numérotée de 1 à 39 dans la marge pour faciliter la correspondance avec les photos des pages du manuscrit.

Le document paraît être une copiée destinée aux archives de la "commune" d'Hauteluce, et contient deux parties: le document d'affranchissement générale concernant les quatre "communes" du mandement de Beaufort ainsi que la procuration passée par les communier d'Hauteluce dans le but de cet affranchissement.

Contexte historique : Le roi de Piemont-Sardaigne et duc de Savoie, Charles-Emmanuel III de Savoie décide de permettre le rachat des droits féodaux par les communautés d’habitants par deux édits : l’édit du 20 janvier 1762, puis celui du 19 décembre 1771 

\textit{"Le texte de 1762 supprimait purement et simplement la taillabilité personnelle envers la
couronne. Concernant les seigneurs privés, elle permettait aux communautés d’habitants de
réaliser un affranchissement collectif, en les chargeant elles-mêmes – via leur syndic - de la
perception d’une indemnité seigneuriale due par chaque personne mainmortable. Après
publication de listes des taillables, les négociations devaient s’engager avec les seigneurs
concernés, sous l’arbitrage des intendants provinciaux. Ces derniers ne purent cependant
contenir toutes les tensions entres habitants et seigneurs, tant au sujet de la validité des listes
qu’au sujet des indemnités et de leur financement. Aussi le texte ne concernait que la taillabilité
personnelle, et c’est là peut-être la raison de sa relative inefficacité. Nombre de charges
seigneuriales relevaient en effet de la taillabilité réelle, celle associée aux terres taillables,
indépendamment de la condition de leurs propriétaires. A quoi bon alors payer pour s’affranchir
d’une taillabilité personnelle si l’on est toujours taillable sur certaines terres ?
Neuf ans plus tard, Charles-Emmanuel III légifère sur un affranchissement général de la
taillabilité personnelle, et réelle : c’est l’édit du 19 décembre 1771. Les affranchissements sont
toujours l’aboutissement de longues négociations qui ne vont pas sans heurts, d’autant plus
qu’on recourt à la vente de biens communaux pour boucler le financement de certaines
procédures. Au final, le résultat est néanmoins significatif. Entre 17541 et 1792, 3 011 contrats
d’affranchissements collectifs furent passés, atteignant la somme de 7 935 409 livres. La très
grande majorité des actes (2 816 sur 3 011) sont postérieurs à l’édit de 17712. Ainsi, on estime
à 65\% le nombre de fiefs affranchis à la veille de la Révolution, cette dernière ayant mis fin à
toute condition servile."} \footnote{\cite{Gachet2022}}

Les communautés du mandement de Beaufort n'ont pas attendu le deuxième décret, et ont finalisé le rachat des droits féodaux dés 1770.

Remarques:
\\- certains feuillets (pages impaires) sont numérotés, il est alors inscrit (Feuillet) en face du numéro
\\- certains une série de paragraphes sont numérotés (1 à 19) mais une autre numérotation (inscrit en début de paragraphe) existe aussi, sans logique apparente
\\- Les mots précédés d'un * sont définis dans un glossaire à la fin du document
\\- il n'a pas été établi d'index des noms de lieux, étant donné que ceux-ci sont peu nombreux et sont répétés dans tout le texte:
les quatre communautés du mandement de Beaufort à savoir :
\begin{list}{}{\leftmargin 1cm}
\item -Saint-Maxime ou St Maxime de Beaufort ou Beaufort
\item -Hauteluce
\item -Villard ou Villard de Beaufort
\item - Queige
\end{list}
Sont cités aussi :
\begin{list}{}{\leftmargin 1cm}
\item -Chambéry , siège de l'administration du duché de Savoie
\item -Turin, capital du Royaume de Piedmont-Sardaigne et du duché de Savoie
\item -la Tarantaise et le Faucigny, deux mandements voisins de celui de Beaufort
\end{list}
\section{Transcription}

\marginnote{\scriptsize Page 1}

Tampon des Archives Communales -Savoie – Propriété Publique
\begin{flushright}
Affranchissement Général \footnote{Pour une étude générale sur les affranchissements voir : \cite{Bruchet}}
\\vente et cession pour les
\\Communautés de St Maxime,
\\Hauteluce, du villard et de Queige
\\mandement de Beaufort
\\ par Le Sgr, marquis du dit lieu 
\\£.81000
\end{flushright}


L'an mille sept cent septante et le seize du mois
de janvier, à Chambéry avant midi dans la
maison du seigneur de Lannoix \index[nom]{Lannoix@\textbf{Lannoix} (Seigneur de) } , par devant
moi, notaire royal soussigné, et présent les témoins 
enfin,  nommées, sest  personnellement établis
et constitués le sieur Pierre Martin contrôleur de 
trésorerie  générale de Savoie\footnote{"le contrôleur de 
trésorerie  générale de Savoie était à l'époque..."}, \index[nom]{Martin@\textbf{Martin} (Pierre) !contrôleur de 
trésorerie  générale de Savoie} natif et habitant
de cette ville lequel degré en qualité de procureur
général et spécial de S.(on) E.(xcellence) messire Joseph 
François Louis Vicardel  Marquis de Fleury \index[nom]{Fleury@\textbf{Fleury}  (Joseph 
François Louis, Vicardel Marquis de) }
ministre d'État, chevalier d'honneur de son
altesse royale madame la Duchesse de Savoie

\marginnote{\scriptsize Page 2}

et Chevalier Grand Croix des ordres de St Maurice 
et Lazare, natif et habitant de la ville de Turin
fils à feu messire François Joseph Nicolas Eléazar
Vicardel Marquis de Fleury, \index[nom]{Fleury@\textbf{Fleury}  (Joseph Nicolas Eléazar, Vicardel Marquis de) }
ministre du cabinet de sa majesté le Roy de 
Pologne, Auguste premier et chevalier de l’ordre 
suprême de Pologne dit de l’Aigle Blanc, par 
procuration du vingt-quatre novembre proche 
passé, reçue par le notaire Ostano \index[nom]{Ostano@\textbf{Ostano}  (Notaire de Turin) }de la ville de 
Turin, dont l’expédition sera jointe à la Minute 
du présent, pour y avoir au besoin recours, et 
relativement aux conventions sous seing privé
faites sous la réserve de l’approbation de sa Majesté 
et non autrement, entre la dite excellence le Seigneur 
Marquis de Fleury\index[nom]{Fleury@\textbf{Fleury}  (Joseph 
François Louis, Vicardel Marquis de) }, et le Sr Michel, fils du Sr 
Joseph Blanc \index[nom]{Blanc@\textbf{Blanc} (Michel, notaire et secrétaire de la communauté de St Maxime) }en qualité de procureur député 
de la communauté de St Maxime de Beaufort, 
par acte du quinze novembre mille sept cent 
soixante-sept Chevallier Joly fils notaire \index[nom]{Chevallier Joly fils@\textbf{Chevallier Joly fils}  (notaire) }, le Sr Joseph 
fils d’autre Joseph Bal notaire \index[nom]{Bal@\textbf{Bal}  (Joseph, notaire et secrétaire de la communauté d’Hauteluce) } et secrétaire de la 


\marginnote{\scriptsize Page 3}

communauté d’Hauteluce a procuré d’icelle
par mandat du dix décembre mille sept cent 
soixante-sept Blanc notaire le Sr Jean-Michel, fils de 
feu Louis Chamiot Maitral \index[nom]{Chamiot Maitral@\textbf{Chamiot Maitral}  (Louis, député de la communauté du Villard et des particuliers de Queige) } député de la 
communauté du Villard et des particuliers de 
Queige, par procuration du vingt décembre même 
année Maitral notaire  des dites conventions en date du 
dix-neuf août mille sept cent soixante-huit, 
sous la même réserve, a affranchi, cédé et vendu, 
ainsi qu’il affranchit, cède et vend par le présent 
en faveur de tous les particuliers des dites communautés 
et paroisses, et aux dites communautés, à l’acceptation 
du dit notaire Michel, fils du dit Sr Joseph Blanc
secrétaire insinuateur\footnote{"un secrétaire insinuateur était à l'époque un notaire en charge d'un bureau d'insinuation (ou du tabellion), c'est-à-dire d'un bureau d'enregistrement par transcription de certains actes civils. Dans les États de Savoie à l'époque Moderne, le terme de tabellion désigne l'ensemble des actes insinués et, par extension, l'administration chargée de la transcription et de la conservation de ces actes. Le bureau du tabellion devait avoir une pièce pour recevoir le public et une pièce pour conserver les actes..."} \index[nom]{Blanc@\textbf{Blanc}  (Joseph, secrétaire insinuateur) }, natif et habitant de la 
paroisse du dit St Maxime de Beaufort, de laquelle 
il est secrétaire et député pour le fait du présent
avec le dit Sr Blanc son père, par la procuration 
cy dessus désignée, et des dits Sr joseph Bal, et
Jean Michel Chamiot-Maitral procureurs députés 
par les dites communautés d’Hauteluce, du Villard 
de Beaufort et des particuliers de Queige à former


\marginnote{\scriptsize Page 4}

des mandats ainsi ci-dessus désignés, ici présents
et acceptants pour les dites communautés, et pour tous
 les dits particuliers et communautés d’ici scavoir,
 il a affranchi et affranchit généralement et
sans aucune exception les dites communautés et les 
particuliers des dites paroisses et communautés et les
particuliers de Queige, de tous devoirs 
seigneuriaux que la dite excellence perçoit ou peut
percevoir \gls{rière} des dites paroisses, en quoi ils consistent
et puissent consister, tant en taillabilité réelle, que
personnelle, comme en *\gls{cens} et censes, \gls{servis} , \gls{laods}, 
\gls{laodsindemnité}, \gls{suffertes}, \gls{dîmes}, \gls{décimes}, 
\gls{corvées}, \gls{alpéages}, \gls{droits de haut siège}, ensemble 
tous hommages, réels et personnels et autres droits 
de la susdite nature, tant à charges honoraires que 
lucratives, tels qu’ils lui appartiennent tant en 
vertu des terriers et reconnaissances que de la dite 
inféodation du vingt-deux décembre mille six 
cent soixante-deux, et des autres acquisitions 
énoncées dans les dites conventions, en sorte que 
dorénavant les biens et personnes des dites

\marginnote{\scriptsize Page 5}

\begin{flushright}
(Feuillet) 3
\end{flushright}	
paroisses soient entièrement libres des dits devoirs 
seigneuriaux, et féodaux et que les mêmes devoirs et
 droits soient annulés, à perpétuité en faveur des 
dites communautés et particuliers, et pour la plus facile 
exécution d’un tel affranchissement général, la dite 
Excellence a cédé et cède aux dites communautés 
respectives le droit d’exiger des \gls{favettiers} les 
sommes auxquelles, en correspectivité de leur 
respectif affranchissement, seront taxés par 
des plans, d'\gls{égances} qui seront dressés, et
 autorisés par la personne que Sa Majesté \index[nom]{Charles-Emmanuel III de Savoie@\textbf{Charles-Emmanuel III de Savoie}  (roi de Sardaigne, duc de Savoie et prince de Piémont)!S.M.} voudra 
commettre et déléguer à ces fins, en outre, la 
dite Excellence\index[nom]{Fleury@\textbf{Fleury}  (Joseph François Louis, Vicardel Marquis de) } a vendu, cédé ainsi que par le 
présent elle vend et cède aux dites communautés 
les \gls{péages}, \gls{bans champêtres}, et tous biens 
immeubles et ruraux qu’elle possède, à lui 
appartenant dans le dit mandement, de la 
possession desquels il s’est dévêtu et en a invêtu 
les susnommés par la tradition d’une plume à 
la manière accoutumée, leur cédant tous noms 

\marginnote{\scriptsize Page 6}

et actions résultant des droits tous interdits, 
\gls{remèdes possessoriaux}, \gls{bénéfices d'instances}, tous
\gls{rescindants}, \gls{rescisoire}, le \gls{droit de pignoratitie} et 
autres semblables, pour la recherche, poursuite et 
exaction desquels le dit Sr. martin \index[nom]{Martin@\textbf{Martin} (Pierre) !procureur du marquis de Fleury} en sa qualité
constitué et substitué pour procureur du seigneur
 marquis et des particuliers, communiers et les dites
communautés en la personne de leurs députés ci-dessus
nommés, et de ceux qu'elles pourront substituer, avec
pouvoir d'en substituer encore d'autres, le tout sous 
élection de domicile, sous la réserve expresse du 
titre de marquis de Beaufort de la juridiction, 
avec tous droits tant lucratifs, honorifiques, qu'autres 
qui y sont relatifs, dans lesquels se trouvent 
compris, la \gls{nomination des officiers de justice}, le 
\gls{greffe}, la \gls{leyde sur le bétail}, la \gls{hale} qu’il possède 
en indivision avec la \gls{mense archiépiscopale} de 
Tarentaise, la chasse, la pêche, le \gls{droit des langues 
dû par les bouchers}, et la \gls{recherche des 
cristaux}, le \gls{droit d’avoir un banc placé} ainsi qu'il

\marginnote{\scriptsize Page 7}

\begin{flushright}
(Feuillet) 4
\end{flushright}	
convient à la qualité des seigneurs du lieu, dans
l'église paroissiale du bourg de St Maxime qui est
la principale terre du dit mandement de Beaufort,
comme encore le droit de faire \gls{placer des
patibulaires avec ses armoiries}, tous lesquels droits
sont ici expressément réservés conformément
aux dites conventions, et desquelles ladite excellence
devra jouir sans la moindre contestation et ses
successeurs à perpétuité. Sans que lesdites Communautés
puissent apporter aucun trouble à cette paisible
jouissance sous quel prétexte que ce soit, pas même
parce que quelques-uns des dits droits ne se
trouveraient pas spécifiés dans ladite inféodation. 
Du vingt-deux décembre mil six cent soixante-
deux, à quoi lesdits députés, en leur qualité
renoncent expressément à l'acceptation du dit Sr
Martin\index[nom]{Martin@\textbf{Martin} (Pierre) !}
, aussi en la qualité, lesdits droits réservés
comme dessus, étant entrés en considération
dans le prix du présent affranchissement auquel
sa dite excellence ne se serait pas déterminée sans

\marginnote{\scriptsize Page 8}

cela, les autres droits ci-dessus vendus et cédés
avec les dits fiefs, se trouvant tant dans le dit
mandement qu’ ailleurs où le tout peut s’étendre
sur les personnes et sur les biens de quelque manière
que ce soit, pour que les dits particuliers, communiers
et communautés jouissent sans autre des droits
ci-dessus cédés et affranchis ainsi qu’elles ont
commencé à jouir dès le premier janvier de
l’année mille sept cent soixante-neuf, relativement
aux conventions privées ci-dessus désignées,
maintenant le dit sieur Martin\index[nom]{Martin@\textbf{Martin} (Pierre) !} en sa qualité, les
droits ci-dessus cédés exempts de tous troubles,
molestier, et autres empêchements de tout le passé
jusqu'à ce jour, se soumettent pour ce à toutes
évictions et garantie a forme du droit, avec
promesse de communiquer aux dits particuliers,
communier, et communautés, à leur première
réquisition moyennant chargé tous les titres,
terriers, extraits, cartes, et autres documents qui
leurs seront nécessaires pour les \gls{égances} et

\marginnote{\scriptsize Page 9}

\begin{flushright}
(Feuillet) 5
\end{flushright}	
répartitions à faire, et généralement tous ceux qui
concernent tous les droits cédés lesquels il promet
leur remettre lorsqu'ils sera entièrement payé, à
l'exception néanmoins de ceux qui affecteront les
droits réservés, dont ils promet seulement la
communication en cas de besoin, pour que les
communautés en puissent prendre des extraits
en forme, promettant les Sr. députés, de leur côté,
que si après la rémission desdits titres et terriers,
le seigneur marquis se trouve en avoir besoin,
pour que ses fermiers puissent exiger les arrérages
qui leur seront dus, ils en feront faire la
communication moyennant charge, étant bien
exprimé que le présent ne pourra préjudicier aux
fermiers et droit ayant du dit seigneur marquis,
tant pour les \gls{servis} engagées que pour les \gls{échuttes},
de tout le passé, jusqu’au premier janvier mille
sept cent soixante-neuf, ce qui sera fait inventaire
de tous les titres, \gls{terriers}, et autres \gls{litterés} dont
le dit seigneur marquis devra faire la rémission

\marginnote{\scriptsize Page 10}

lors de son payement et dont il a ci de plus promit
la communication. Le présent affranchissements
général, cession et vente faits pour le prix de quatre
vingt mille livres, et mille livres pour épingles,
lesquelles mille livres ont été présentement et
réellement comptées, savoir cinq cent trente-
huit livres dix-huit sols dix deniers par ledit
Sr Blanc en sa qualité, en vingt-deux pistoles
neuves, valant chacune vingt-quatre livres et le
surplus en sols et pièces de deux deniers, trois cent
vingt-quatre livres neuf sols deux deniers par le
dit Sr Bal aussi sa qualité en treize pistoles
et demi, et le surplus en sols et une pièce de deux
deniers, et enfin trente-six livres douze sols par
le dit Sr Maitral  en sa qualité en cinq pistoles
et demi et le surplus en sols le tout par ledit Sr .
Martin\index[nom]{Martin@\textbf{Martin} (Pierre) !} vérifié, retiré et emboursé au vu de moi
notaire et témoins, dont il se contente et quitte,
lesquels payements ils ont fait par proportion à la
cote de chacune des communautés relativement

\marginnote{\scriptsize Page 11}

\begin{flushright}
(Feuillet) 6
\end{flushright}	
aux conventions faites entre elles, le cinq juillet
mil sept cent soixante-huit qui seront ténorisées
enfin du présent avec celles faites avec le dit
seigneur marquis qui feront corps à icelui. Les
parties déclarant s'y rapporter pour tout ce qui
ne se trouvera pas suffisamment expliqué, eu
égard que les payements de ladite somme de quatre-
vingt mille livres devra se faire suivant les dites
conventions, et laquelle susdite somme, les dits
communiers et communautés promettent payer
solidairement par le ministère de leurs députés,
qui à ces fins renoncent à tous bénéfices de
division, d’ordre, et de discussion dans le terme
de dix ans, et en différents payements dont
chacun ne sera pas moindre de huit mille livres,
en avertissant néanmoins le dit seigneur marquis
deux mois auparavant, lequel seigneur sera obligé
de faire recevoir les dits payements en Savoie,
dans la présente ville par le ministère d'un
procureur qui sera à ces fins constitué, et auquel
on pourra faire les avertissements ci-dessus stipulés

\marginnote{\scriptsize Page 12}

et jusqu’au dit paiement les susnommés toujours 
solidairement et  avec les mêmes renonciations,
promettent payer l’intérêt de ladite somme
totale de quatre-vingt mille livres aussi dans
cette ville entre les mains de celui qui sera
député par le dit seigneur marquis, et c’est à raison
de trois pour cent, lequel intérêt sera diminué
à proportion des paiements qui seront faits de
la manière ci-dessus convenue, et sera payé à
l’échéance de chaque année relativement
aux susdites conventions, suivant lesquelles il doit
être payé de la dite somme totale celle de trente-
six mille sept cent cinq livres quatorze sols
par la communauté de Ste Maxime, trente-deux
mille trois cent soixante-six livres six sols
par celle d’Hauteluce, et dix mille neuf cent
vingt-huit livres par celles du Villard et  Queige,
sans que la division faite entre les dites communautés
pour le paiement de la susdite somme totale
de quatre-vingt mille livres puisse en aucune
manière préjudicier à l’obligation solidaire par
elles ci-devant contractée en faveur de Sa dite

\marginnote{\scriptsize Page 13}

\begin{flushright}
(Feuillet) 7
\end{flushright}	
excellence le seigneur marquis de Fleury, entre
toutes lesquelles communautés les immeubles ci-
dessus vendus, restent indivis, avec convention que
sur les paiements qui seront faits en déduction de
ladite somme capitale, les premiers deniers seront
employés au paiement de la mieux value réservée
au royal patrimoine, le seigneur marquis restant
chargé d’en faire faire l’acceptation de la
manière qui plaira à Sa Majesté de l’ordonner.
Le dit seigneur se chargeant également de fournir
une application assurée, et de faire à ces fins toutes
procédures nécessaires pour regard de la somme
qui sera déterminée devant être appliquée pour
raison des fidéicommis et primogénitures qui se
trouvent dans la famille dudit seigneur marquis,
laquelle application devra être faite suivant
les déterminations qui seront données par
Sa Majesté sur les instances qui seront faites de
la part dudit seigneur marquis, et tout ce que
de plus a été promis observer, tant par le dit
Sieur. Martin\index[nom]{Martin@\textbf{Martin} (Pierre) !}
que par les députés ci-dessus nommés

\marginnote{\scriptsize Page 14} 

avec promesse de le faire observer chacun, en
ce qui les concerne, aux mêmes respectives de tous
dépends \gls{doages}, intérêts et à l’obligation et
constitution, savoir de la part du dit Sieur Martin\index[nom]{Martin@\textbf{Martin} (Pierre) !},
des biens du dit seigneur marquis de Fleury,
qui se constitue tenir en vertu de ladite
procuration, et de la part des dits députés
ceux des dits communiers et des dits communautés,
qui se constituent également tenir en vertu
des procurations ci-dessus énoncées, les expéditions
desquelles seront dûment jointes à la minute
du présent avec l’approbation faite par la
communauté du village de Beaufort, par
acte du vingt-huit décembre mil sept cent
soixante-huit. Sieur Maitral notaire \index[nom]{Maitral@\textbf{Maitral} (notaire)}fait et prononcé
au dit lieu en présence du sieur François Chevallier,
marchand \index[nom]{Chevallier@\textbf{Chevallier} (François, marchand, témoin)}, et de maître Pierre Dominique Richard \index[nom]{Richard@\textbf{Richard} (Dominique, substitut procureur du Sénat, témoin)},
substitut procureur du Sénat, tous deux habitants
de la présente ville, témoins requis.

\marginnote{\scriptsize Page 15}

\begin{flushright}
(Feuillet) 8
\end{flushright}

\begin{flushright}
Teneur de la procuration passée		
\\par les communiers d’Hauteluce		
\\à				
\\Maitre Joseph Bal, notaire secrétaire		
\\du dit Hauteluce\index[nom]{Bal@\textbf{Bal}  (Joseph, notaire et secrétaire de la communauté d’Hauteluce) }.	
\end{flushright}

L’an mille sept cent soixante-sept, et le sixième
jour du mois de décembre, à l’issue de la messe
paroissiale d’Hauteluce, et au \gls{banc du droit}
du dit lieu, ainsi qu’au conseil de la dite paroisse
d’Hauteluce, avec celles de St. Maxime et du Villard
qui, avec la paroisse de Queige, composent
le territoire et marquisat de Beaufort, avaient
dés longtemps projetés, à l’exemple des paroisses
leurs voisines de Faucigny et de Tarentaise,
d’acquérir les fiefs qui dépendent du dit marquisat, et
tous autres fiefs qui peuvent s’étendre
sur le territoire des dites paroisses, aux fins de
s’affranchir de toute servitude et de rendre
libres pour leurs personnes, ainsi que leurs biens de
tout hommages personnels et réels d’en supprimer
et anéantir les effets, et les motifs en étaient justes

\marginnote{\scriptsize Page 16}

et louables, puisqu’il en devait résulter un grand
avantage pour les particuliers des dites communautés,
et les conseils des dites trois paroisses appuyées du
suffrage du plus grand nombre et des plus
considérables habitants d’icelles auxquels ils en
avaient fait part, en firent faire leur  proposition
respectueuse au seigneur marquis de Fleury et
de Beaufort\index[nom]{Fleury@\textbf{Fleury}  (Joseph François Louis, Vicardel Marquis de) !marquis de Beaufort}, qui par une suite de ses bontés
ordinaires pour les dites communautés leur fit
savoir qu’il ne désapprouvait point leur
démarche, et donnerait volontiers les mains,
leur envoyant à cet effet une note des différents
fiefs qu’il avait dans le mandement du dit Beaufort
et du prix des acquisitions qui en avaient été faites,
laquelle est ci-après ténorisée.

\begin{center}
2
\end{center}
1° La Baronnie de Beaufort érigée ensuite en
marquisat en faveur du Seigneur marquis Joseph
François Vicardel de Fleury\index[nom]{Fleury@\textbf{Fleury}  (Joseph Nicolas Eléazar, Vicardel Marquis de) }, l’un de ses autheurs,
a été échangée par contrat du vingt-deux décembre
mille sept cent soixante-deux Marthar notaire\index[nom]{Marthar@\textbf{Marthar} (notaire)}


\marginnote{\scriptsize Page 17}

\begin{flushright}
(Feuillet) 9
\end{flushright}		
contre le palais des ambassadeurs dont le dit
seigneur marquis avait peu auparavant fait
l’acquisition du marquis Vecuso de la Tour\index[nom]{Vecuso de la Tour@\textbf{Vecuso de la Tour}  (Marquis) },
pour le prix de seize mille cinq cent ducatons
de la valeur pour lors de quatre livres quinze
sols pièce, prix auquel il a été estimé par ledit
contrat d’échange, et à la dite terre de Beaufort
à cette somme évaluée, ainsi qu’à celle de
deux mille sept cent cinquante florins, de
mieux value, soit nonante cinq florins de
rente annuelle suivant la patente d’inféodation
donnée par la chambre des comptes, par laquelle
patente d’inféodation le droit de rachat
perpétuel est réservé à la couronne lesquels
seize mille cinq cent ducatons sur le pied
porté par ledit état reviennent à septante 
huit mille trois cent septante cinq livres,

\begin{center}
3
\end{center}
2°La rente de Gerbais, acquise ensuite par le
même seigneur de Fleury, du Seigneur baron
de Blancheville \index[nom]{Blancheville@\textbf{Blancheville}  (baron de) }
 par contrat du vingt-huit

\marginnote{\scriptsize Page 18}

décembre mille six cent soixante-quatre, Dardel 
notaire 
\index[nom]{Dardel@\textbf{Dardel} (notaire)}
, rapportant mille livres florins de revenu 
pour le prix de vingt-trois mille huit cent florins, 
revient à raison de treize sols quatre deniers 
chaque florin, soit les trois florins faisant deux 
livres, à quinze mille huit cent soixante-six 
livres treize sols quatre deniers.

\begin{center}
4
\end{center}
3°La rente de Bellegarde acquise par le Seigneur
 Nicolas Éléazar Vicardel, marquis de Toivier\index[nom]{Fleury@\textbf{Fleury}  (Joseph Nicolas Eléazar, Vicardel Marquis de) !marquis de Toivier}, des
 messieurs d’Ancieu\index[nom]{Ancieu@\textbf{Ancieu} (messieurs d’)}, pour le prix de mille florins, et
 la rente, d’Arnolled autre petit fief qu’il a acquis
 par acte du vingt-quatre décembre septembre
 mil six cent nonante-huit, Bertier Notaire \index[nom]{Bertier@\textbf{Bertier}  (notaire)} pour le prix
 de huit cent florins, lesquels deux pris
sur le même pied que ci-dessus arrivent à 
deux mille huit cent livres.

\begin{center}
5
\end{center}
4° Le moderne marquis a de plus acquis par contrat
 du vingt-sept août mil sept cent quarante-quatre
reçu par moi notaire de Joseph Vibert \index[nom]{Vibert@\textbf{Vibert}  (Joseph, notaire)}un petit fief

\marginnote{\scriptsize Page 19}

\begin{flushright}
(Feuillet) 10
\end{flushright}		
appelée la Rente d’Anvey ( ?) pour le prix de mille
 deux cent livres.

\begin{center}
6
\end{center}
Toutes lesquelles sommes réunies emportent la
 totale de nonante-six mille six cent quarante-une 
livres, treize sols, quatre deniers, de laquelle suivant 
la lettre du seigneur marquis, prélevant pour le 
titre, droit de juridiction, et environ pour trois 
cent livres de droits seigneuriaux qu’il dit en 
être inséparables, il resterait donc quatre-vingt 
mille livres en capital principal, ou deux mille 
quatre cents livres de rente annuelle au trois pour 
cent qu’il paye aux communautés.

\begin{center}
7
\end{center}
D’après la réception duquel état donné, pour
ledit seigneur marquis, les conseils des dites trois 
paroisses s’étant assemblés, ils auraient pour 
parvenir au but dudit affranchissement délibérés, 
ainsi qu’ils l’ont fait par l’acte de délibération 
motivé du onze octobre dernier, de supplier le 
seigneur intendant général de permettre dans chacune 
des dites trois paroisses la convocation d’une

\marginnote{\scriptsize Page 20}

assemblée générale, et elle leur fut permise par
décret mis en fin d’icelle délibération du vingt-
neuf même mois, avec commission au premier notaire
sur ce requis d’en dresser par les actes et verbaux,
en conséquence des syndics et conseil de cette Communauté
d’Hauteluce, prièrent le Révérend Curé de faire les
annonces de la dite assemblée générale, dimanche
dernier vingt-neuf novembre à son prône, ainsi
qu’il se pratique en tel cas, ce qui fut exécuté, et
le dit jour à l’issue de la messe paroissiale, le
Sergent Vionnet \index[nom]{Vionnet@\textbf{Vionnet}  (Sergent) } en fit au présent banc du droit
de nouveau l’avertissements avec la lecture de la
dite délibération et décret dont il afficha une copie
au dit banc du droit, avec intimation qu’il fit à
tous les particuliers, chefs de famille de cette
paroisse, de s’assembler à ces fins ce jourd’hui à
la présente heure a formé de son exploit du dit jour,
et outre le peuple en à de nouveau jour ce jourd’hui été
instruit et convoqué tant par le son de la cloche
à la manière accoutumée, par de nouvelles annonces


\marginnote{\scriptsize Page 21}

\begin{flushright}
(Feuillet) 11
\end{flushright}
au \gls{prône}, et par de nouvelles criées de la part du
dit Sergent Vionnet\index[nom]{Vionnet@\textbf{Vionnet}  (Sergent) }, à l’issue de la messe paroissiale,
par lesquelles il a averti le peuple en la personne des
chefs de famille de s’assembler au présent lieu du
banc du droit pour délibérer sur le dit affranchissement,
et de la manière de l’exécuter, lequel a été exécuté
comme s’ensuit

\begin{center}
8
\end{center}
À ces causes par devant, moi Michel Blanc, notaire 
et secrétaire de la communauté de St. Maxime\index[nom]{Blanc@\textbf{Blanc} (Michel, notaire et secrétaire de la communauté de St Maxime) },
soussigné, et en présence des témoins enfin nommés, se
sont personnellement établis et constitués les susnommés
syndic, conseillers et communiers de la paroisse
d’Hauteluce, savoir : 
Rev. Sr. François Girod Roux prieur et curé d’icelle, \index[nom]{Girod Roux (François, Rev. Sr. , Prieur et curé d'Hauteluce) }
Rev. Sr. Claude Baptandier prêtre et vicaire, \index[nom]{Baptandier (Claude, Rev. Sr. , prêtre et vicaire) }
Rev. Sr. Pierre Reymond prêtre et régent, \index[nom]{Reymond (Pierre, Rev. Sr. , prêtre et régent) }
Honorable Joseph Baptandier syndic, \index[nom]{Baptandier (Joseph, Honorable, Syndic) }
Joseph Bal Fourrer conseiller, \index[nom]{Bal Fourier (Joseph, Conseiller) }
Barthélémy Vibert  Lechet, \index[nom]{Vibert Lechet (Barthélémy) }
Donaz Vionnet, \index[nom]{Vionnet (Donaz) }
Louis Vionnet Fuasset, \index[nom]{Vionnet Fuasset (Louis) }
et, 
Nicolas Desduc, \index[nom]{Desduc (Nicolas) }
tous membres du dit conseil,
\\
\marginnote{\scriptsize Page 22}

Louis Bal Fontaine, \index[nom]{Bal Fontaine (Louis) }
Pierre Mailliand Gonod pour son père malade, \index[nom]{Mailliand Gonod (Pierre, pour son père malade) }
Nicolas Thovex Comtoz dit la Salaz, \index[nom]{Thovex Comtoz (Nicolas, dit Salaz) }
Jacques Rey, \index[nom]{Rey (Jacques) }
Nicolas fils de Donat Thovex Comtoz, \index[nom]{Thovex Comtoz (Nicolas, fils de Donat Thovex Comtoz) }
Joseph Bal Pétré, \index[nom]{Bal Pétré (Joseph) }
Michel à feu Pierre Guiguet, \index[nom]{Guiguet (Michel, à feu Pierre Guiguet) }
Joseph Piccard, \index[nom]{Piccard (Joseph) }
François à feu Louis Favre, \index[nom]{Favre (François, à feu Louis Favre) }
Jacques Pichol, \index[nom]{Pichol (Jacques) }
Pierre Ducretel, \index[nom]{Ducretel (Pierre) }
Philippe Vionnet dit Gabiollet, \index[nom]{Vionnet (Philippe) }
François Perrier, \index[nom]{Perrier (François) }
Jacques à feu Joseph Pichol, \index[nom]{Pichol (Jacques, à feu Joseph Pichol) }
Louis Braisaz, \index[nom]{Braisaz (Louis) }
Claude Guiguoz, \index[nom]{Guiguoz (Claude) }
André Lachenal Siordet, \index[nom]{Lachenal Siordet (André) }
Pierre à feu Jean Baptiste Bal Pétré, \index[nom]{Bal Pétré (Pierre, à feu Jean Baptiste Bal Pétré) }
Louis à feu Jacques Guiguet, \index[nom]{Guiguet (Louis, à feu Jacque Guiguet) }
Joseph Alex, \index[nom]{Alex (Joseph) }
Pierre Ligeon, \index[nom]{Ligeon (Pierre) }
Joseph Constantin, \index[nom]{Constantin (Joseph) }
Jean-Pierre Bardet Bardillion, \index[nom]{Bardet Bardillion (Jean-Pierre) }
Jacques Palluel Blanc, \index[nom]{Palluel Blanc (Jacques) }
Jean Roset, \index[nom]{Roset (Jean) }
Joseph Franc Carmetrant, \index[nom]{Franc Carmetrant (Joseph) }
Claude Bal, \index[nom]{Bal (Claude) }
Michel à feu Jean Guiguet, \index[nom]{Guiguet (Michel, à feu Jean Guiguet) }
Pierre Duret Viallet, \index[nom]{Duret Viallet (Pierre) }
André Braisaz Cretet, \index[nom]{Braisaz Cretet (André) }
Jacques Bouchage, \index[nom]{Bouchage (Jacques) }
Mammert Deduc, \index[nom]{Deduc (Mammert) }
Joseph Cuvex Combaz, \index[nom]{Cuvex Combaz (Joseph) }
Pierre Verpil, \index[nom]{Verpil (Pierre) }
Louis Bron Bozon, \index[nom]{Bron Bozon (Louis) }
Alexis Garzent, \index[nom]{Garzend (Alexis) }
Sr. Jean-Baptiste Christine, \index[nom]{Christiné (Jean-Baptiste) }
François Guiguoz Cuvex, \index[nom]{Guiguoz Cuvex (François) }
Jean-Baptiste Guigoz, \index[nom]{Guigoz (Jean-Baptiste) }
Pierre Bochet Fenet, \index[nom]{Bochet Fenet (Pierre) }
Claude Mailliand, \index[nom]{Mailliand (Claude) }
Michel Pichol, \index[nom]{Pichol (Michel) }
Jacques Mailliand, \index[nom]{Mailliand (Jacques) }
\\
\marginnote{\scriptsize Page 23}

Jean-Baptiste Durand, \index[nom]{Durand (Jean-Baptiste) }
Jean-Baptiste Bardet Bardillon, \index[nom]{Bardet Bardillon (Jean-Baptiste) }
Jean François Bal Pétré, \index[nom]{Bal Pétré (Jean François) }
Nicolas Mailliand Gonod, \index[nom]{Mailliand Gonod (Nicolas) }
Joseph Braisaz, \index[nom]{Braisaz (Joseph) }
Nicolas Cuvex Micholin, \index[nom]{Cuvex Micholin (Nicolas) }
Joseph Guiguoz, \index[nom]{Guiguoz (Joseph) }
Jean-Jacques Vionnet Ducrey, \index[nom]{Vionnet Ducray (Jean-Jacques) }
Pierre Blanchard, \index[nom]{Blanchard (Pierre) }
Claude Christian, \index[nom]{Christian (Claude) }
François Martin, \index[nom]{Martin (François) }
Claude Durand, \index[nom]{Durand (Claude) }
Antoine Avocat, \index[nom]{Avocat (Antoine) }
Michel Baptamdier Santiquet, \index[nom]{Baptandier Santiquet (Michel) }
François fils de Pierre Pichol pour son père malade, \index[nom]{Pichol (François, fils de Pierre Pichol, pour son père malade) }
Joseph Poupelloz, \index[nom]{Poupelloz (Joseph) }
François Verpil, \index[nom]{Verpil (François) }
Aymé Cuvex Combaz, \index[nom]{Cuvex Combaz (Aymé) }
Claude Mailliand Gonod, \index[nom]{Mailliand Gonod (Claude) }
Joseph Vionnet Fuapset, \index[nom]{Vionnet Fuapset (Joseph) }
Pierre à feu Jean François Bal, \index[nom]{Bal (Pierre, à feu Jean François Bal) }
Donat Mailliand Gonod, \index[nom]{Mailliand Gonod (Donat) }
Pierre à feu Lachenal Siordet, \index[nom]{Lachenal Siordet (Pierre, à feu Lachenal Siordet) }
Jacques Poupelloz, \index[nom]{Poupelloz (Jacques) }
Joseph Bron Bozon, \index[nom]{Bron Bozon (Joseph) }
Mathieu Bal Pétré, \index[nom]{Bal Pétré (Mathieu) }
Claude-François Babullin, \index[nom]{Babullin (Claude-François) }
Antoine Guiguet, \index[nom]{Guiguet (Antoine) }
Claude Braisaz, \index[nom]{Braisaz (Claude) }
Jean-Pierre Viard, \index[nom]{Viard (Jean-Pierre) }
Mermet Favre Coutillet, \index[nom]{Favre Cotillet (Mermet) }
Mermert à feu André Palluel Blanc, \index[nom]{Palluel Blanc (Mermet, à feu André Palluel Blanc) }
Pierre à feu Nicolas Guiguet, \index[nom]{Guiguet (Pierre, à feu Nicolas Guiguet) }
André Favre Cotillet, \index[nom]{Favre Cotillet (André) }
Le sieur Donat Bal, \index[nom]{Bal (Donat) }
Joseph Roger, \index[nom]{Roger (Joseph) }
François Viardet, \index[nom]{Viardet (François) }
Joseph Perrouz, \index[nom]{Perrouz (Joseph) }
Jacques Bal Pétré, \index[nom]{Bal Pétré (Jacques) }
Jean-Baptiste Vibert Bourgeois, \index[nom]{Vibert Bourgeois (Jean-Baptiste) }
Joseph Favre Cotillet, \index[nom]{Favre Cotillet (Joseph) }
Pierre Cuvex Combaz, \index[nom]{Cuvex Combaz (Pierre) }

\marginnote{\scriptsize Page 24}

Nicolas Bal, \index[nom]{Bal (Nicolas) }
Claude Duc Pithiod, \index[nom]{Duc Pithiod (Claude) }
Mammert à feu Jean Baptiste Palluel Blanc, \index[nom]{Palluel Blanc (Mammert, à feu Jean Baptiste Palluel Blanc) }
Joseph à feu François Bal, \index[nom]{Bal (Joseph, à feu François Bal) }
Claude Avocat, \index[nom]{Avocat (Claude) }
André Nicolas Gamin, \index[nom]{Gamin (André Nicolas) }
André Ligeon, \index[nom]{Ligeon (André) }
Joseph Ducretet, \index[nom]{Ducretet (Joseph) }
Mammert Ducretet, \index[nom]{Ducretet (Mammert) }
Claude Favre Cotillet, \index[nom]{Favre Cotillet (Claude) }
Joseph Ligeon Ligeonnet, \index[nom]{Ligeon Ligeonnet (Joseph) }
Sylvestre Lachenal Viordet, \index[nom]{Lachenal Siordet (Sylvestre) }
Jean-Baptiste Ducretet, \index[nom]{Ducretet (Jean-Baptiste) }
François Carcey, \index[nom]{Carcey (François) }
Joseph-Marie Guiguoz, \index[nom]{Guiguoz (Joseph-Marie) }
Jean-Baptiste Guiguet, \index[nom]{Guiguet (Jean-Baptiste) }
Pierre Ligeon comme tuteur de Donza Christiné, \index[nom]{Ligeon (Pierre, tuteur de Donza Christiné) }
Jean-Baptiste Ligeon Ligeonnet à feu Jacques, \index[nom]{Ligeon Ligeonnet (Jean-Baptiste) }
Jean-Baptiste Ligeon, \index[nom]{Ligeon (Jean-Baptiste) }
Jean-Baptiste Ralard., \index[nom]{Ralard (Jean-Baptiste) }
\\
À tous lesquels se disent excéder les deux tiers des
chefs de famille de ladite paroisse d’Hauteluce,
et y étant actuellement, les trois faisant le tout,
lesquels, après avoir ouï la lecture de la délibération
du douze octobre, et du décret de péréquation du
seigneur intendant général, soit du Sr. sus-délégué,
de la lettre du dit seigneur marquis de Beaufort
et ses propositions, et après avoir fait leurs réflexions
sur leur contenu, et sur ce qui pouvait y avoir du 


\marginnote{\scriptsize Page 25}

rapport, tant à leur nom qu’à celui des autres
absents, ont convenu et délibéré d’une voix
unanime, comme s’ensuit :
\begin{center}
9
\end{center}
1° Que ses députés qui seront élus, feront leur
très humble supplication à S.M.\index[nom]{Charles-Emmanuel III de Savoie@\textbf{Charles-Emmanuel III de Savoie}  (roi de Sardaigne, duc de Savoie et prince de Piémont)!S.M.} pour qu’il lui
plaise leur permettre d’acquérir conjointement avec
les communiers de Beaufort et du Villard du dit
seigneur marquis de Fleury\index[nom]{Fleury@\textbf{Fleury}  (Joseph François Louis, Vicardel Marquis de) }, les fiefs à lui appartenant
et qui dépendent du dit seigneur marquis
de Beaufort, que les autres ci-devant mentionnés
sous la seule réserve en faveur dudit seigneur marquis
du titre de la juridiction et des droits qui y sont
purement relatifs, tant honorifiques qu’autre,
et que pour la validité de telles acquisitions, il plaise
à S.M.\index[nom]{Charles-Emmanuel III de Savoie@\textbf{Charles-Emmanuel III de Savoie}  (roi de Sardaigne, duc de Savoie et prince de Piémont)!S.M.} ou de se départit du \gls{droit de reachapt}
réservé par les patentes d’inféodation, sous
l’offre de payer telle et modérée somme, à la-
quelle ils pourront être taxés. L’intention des dits

\marginnote{\scriptsize Page 26}

communiers étant d’éteindre et d’anéantir pour eux et
leurs successeurs et à perpétuité, tous hommages,
quelconques personnels et réels et de rendre les
personnes et les biens astreints et sujets aux dits fiefs
entièrement libres, en telle sorte qu’il ne reste
par rapport à leurs personnes et à leurs biens, aucun
droit des dits fiefs, en quoi qu’ils consistent et puissent consister
sauf la simple juridiction, \gls{droit de pêche de chasse},
\gls{Laide}, minéraux qui sont les seuls droits qu’elles
prétendent excepter de leur acquisition, tout le surplus
des dits fiefs devant être compris dans l’acquisition
qu’ils se proposent aux fins de l’entière extinction
de tous autres droits perçus ou à percevoir pour raison
des dits fiefs en dépendance d’iceux ou à raison du
dit marquisat, tant riere le territoire des trois paroisses,
que riere la paroisse de Queige, et ailleurs où les dits
fiefs peuvent s’étendre sur les personnes et sur les
biens.

\begin{center}
10
\end{center}
À ces fins ils donnent à leurs députés pouvoir autant
ample qu’il pourra être nécessaire pour leur recours

\marginnote{\scriptsize Page 27}

à S.M., traiter et négocier avec les Royales
Finances tant pour raison des fiefs appartenants au dit
seigneur marquis de Beaufort \index[nom]{Fleury@\textbf{Fleury}  (Joseph François Louis, Vicardel Marquis de) !marquis de Beaufort} qu’à d’autres seigneurs
dans les territoires et appartenances des dites communautés,
et pour qu’il lui plaise leur permettre de faire
l’acquisition des dits fiefs, moyennant finance, donnant
par avance tout plus ample pouvoir à leurs députés
pour le faire des acquisitions au prix qu’il avisera.

\begin{center}
11
\end{center}
2° Ils autorisent leurs députés à se prévaloir à leur
nom des offres et déclarations à eux faites par
le dit seigneur marquis de Beaufort\index[nom]{Fleury@\textbf{Fleury}  (Joseph François Louis, Vicardel Marquis de) !marquis de Beaufort} conjointement
aux communautés de St. Maxime, du Villard, de
traiter avec le dit seigneur au prix de quatre-vingt
mille livres, ils ne peuvent obtenir un meilleur
traitement et condition, à même ils s’en remettent
à sa direction pour les épingles qu’il conviendra de
donner
\begin{center}
12
\end{center}
3° Leur député devra stipuler tant au regard du dit
seigneur marquis, qu’autres seigneurs de fiefs

\marginnote{\scriptsize Page 28}

Que les communiers de leurs paroisses et ceux de Ste-
Maxime et du Villard qui s’unissent pour faire les
dites acquisitions, pourront libérer du prix d’icelles
en divers payements et lorsqu’ils seront en faculté de
le faire, pourvu que cependant quant au dit seigneur
marquis chaque payement ne soit pas moindre de
huit mille livres, et quant aux autres seigneurs de fiefs,
de deux mille livres, lorsque le prix doublera cette
somme, et que jusqu’au dit payement il en sera
payé la cense qui sera convenue par leurs députés au
dernier les moins onéreux pour les communautés.

\begin{center}
13
\end{center}
4° Que dans le cas, qui suivant les négociations il
fallût trouver de l’argent comptant, pour s’acquitter
en partie du prix des acquisitions, leur député pourra
faire les emprunts nécessaires, tant par constitution de
rente qu’autrement, et obliger pour la restitution
des capitaux et payement des censes ou intérêt,
tous les biens de la communauté en général, et ceux

\marginnote{\scriptsize Page 29}

des particuliers qui sont intervenus au présent, présents
et avenir, aux constitutions d’icelles, et pour tous
dépends que pourraient supporter les prêteurs.
\\
\begin{center}
14
\end{center}
5° Ils donnent en outre pouvoir à leur député de faire
tels emprunts qu’ils jugeront être nécessaires, et par
défaut de fonds, pour fournir aux frais indispensables
qu’ils auront à supporter pour lesdites négociations,
frais de transport dans la ville de Turin, séjour
dans icelle, qu’autres même pour payer les sommes
auxquelles ils seront taxés par la finance des dites
acquisitions, à quelles fins ils supplient S.M.\index[nom]{Charles-Emmanuel III de Savoie@\textbf{Charles-Emmanuel III de Savoie}  (roi de Sardaigne, duc de Savoie et prince de Piémont)!S.M.}
de les autoriser pour ce regard.
\\
\begin{center}
15
\end{center}
6°Il devra encore supplier très humblement S.M.\index[nom]{Charles-Emmanuel III de Savoie@\textbf{Charles-Emmanuel III de Savoie}  (roi de Sardaigne, duc de Savoie et prince de Piémont)!S.M.}
de leur permettre la vente et aliénation des dîmes
et autres choses aliénables qui seront comprises dans
les ventes à passer avec ledit seigneur marquis
et autres seigneurs de fiefs, ensemble la vente et 

\marginnote{\scriptsize Page 30}

aliénation des communaux qui à leur paroisse
d’Hauteluce, et en laissant toute fois une quantité
suffisante pour inalper les bestiaux qui s’hivernent
dans l’endroit, pour être les deniers en provenant
employés à payer en partie le prix des acquisitions,
la finance ou les frais à supporter par le moyen de
laquelle vente, les acquisitions des dits fiefs deviendront
moins onéreuses à la communauté, et pour vérifier,
connaitre et fixer les communaux ou portions iceux
qui devront être vendues, ils nomment par avance
trois personnes gens experts, à savoir le Sr. Joseph
à feu François Bal\index[nom]{Bal@\textbf{Bal} (Joseph)}, Honorable Pierre à feu Jean-François
Bal Pétré\index[nom]{Bal Pétré@\textbf{Bal Pétré} (Pierre)} et Claude à feu Claude Braisaz\index[nom]{Braisaz@\textbf{Braisaz} (Claude)}, tous trois
du dit Hauteluce, avec offre d’y faire intervenir
un géomètre député de la part du Sgnr., intendant
général, le syndic ou un conseiller et le secrétaire
de la communauté, lequel secrétaire recevra et
réduira par écrit en forme de verbal tout ce qui

\marginnote{\scriptsize Page 31}


aura été fait, vérifié, arrêté et convenu pour ce
regard, et qu’en autorisant la communauté à vendre
ainsi partie de ses communaux, il plaise à S.M.\index[nom]{Charles-Emmanuel III de Savoie@\textbf{Charles-Emmanuel III de Savoie}  (roi de Sardaigne, duc de Savoie et prince de Piémont)!S.M.}
leur accorder un délégué auquel devra s’adresser
leur député pour pouvoir les faire exposer en vente
et enchères, et auquel il appartiendra d’en approuver et
autoriser les expéditions qui en seront faites aux
enchérisseurs, lesquels soient tenus d’en délivrer le
prix entre les mains du receveur que le conseil
établira, et dont la quittance suffira pour leur
plein payement acquittement.
\\
\begin{center}
16
\end{center}
7°Et comme plusieurs ont usurpés sur les communaux,
surtout les aboutissants, ils donnent pouvoir à
leur dit député de faire recherche des usurpations
et de traiter avec ceux qu’il connaîtra être dans
le dit cas, régler et convenir avec eux du
relâchement auquel ils seront tenus, même à prix
d’argent, et dans le cas que ceux qui auront ainsi

\marginnote{\scriptsize Page 32}
usurpés, ne veuillent régler à l’amiable, leurs députés
en autorisé par le présent à les y contraindre par
les voies judiciaires.
\\
\begin{center}
17
\end{center}
8°Qu’attendu que le prix qui pourra retirer des
parties des communaux ne pourra faire qu’une partie
du montant des dépenses, et pour en venir au moyen
d’apurer, il restera autant des grosses sommes, et pour
qu’un chacun supporte la charge à proportion du bénéfice
qu’il en doit ressentir, il est de nécessité d’en venir à
une répartition, à ces fins leurs députés suppliera
S.M.\index[nom]{Charles-Emmanuel III de Savoie@\textbf{Charles-Emmanuel III de Savoie}  (roi de Sardaigne, duc de Savoie et prince de Piémont)!S.M.} de leur permettre de faire procéder à
une estime des servitudes auxquelles chaque particulier
peut être astreint pour sa personne ou pour ses
biens par des commissaires qui seront choisis à
ces fins par leurs dits députés, laquelle devra être
basse et ménagée, pour que la charge leur soit
rendue plus supportable, à laquelle taxe en fait
de servis ne devra excéder un capital formé sur le
pied du cinq pour cent du montant d’iceux

\marginnote{\scriptsize Page 33}


et de ladite taxe il en sera dressé un état [...]  
a devoir être approuvé par le seigneur intendant général
étant loisible aux dits particuliers de se libérer de ce que
chacun se trouvera taxé entre les mains du même
receveur ci-devant a établir par le conseil et
jusqu’à ce ils  en payeront annuellement l’intérêt
sur le même pied que la communauté sera obligée
de le payer pour le prix des dites acquisitions ou sur
le même pied qu’elle le payera en cas de d’emprunt
en venant à avoir lieu à faire des payements par
remboursement ou autrement pour le bien de la
communauté, par la liberté qu’elle aura de s’acquitter
en partie brisée, et en ce cas les dits particuliers pourront
être contraints d’en payer les capitaux, et le surplus
à quoi pourrait arriver le prix des dites acquisitions
finance frais à faire en conséquence au-delà tant du
produit de la vente des communaux que des taxes
à imposer aux dits particuliers, sur le pied ci-dessus.
Il sera porté sur le bilan de la cotte générique
et rapporté sur la taille d’un chacun, et comme
les fiefs qu’on se propose d’acquérir tant du dit

\marginnote{\scriptsize Page 34}

Seigneur marquis qu’autres peuvent s’étendre pour
quelque portion sur la paroisse de Queige, et autres
terrains du voisinage, S.M.\index[nom]{Charles-Emmanuel III de Savoie@\textbf{Charles-Emmanuel III de Savoie}  (roi de Sardaigne, duc de Savoie et prince de Piémont)!S.M.} sera suppliée qu’on
exige des particuliers qui seront contraints à quelques
servitudes, pour leur personne ou pour leurs biens,
l’estime et la taxe qui en sera faite par des commissaires
nommés respectivement,

\begin{center}
18
\end{center}
9° En conséquence de ce qui a été convenu entre le dit
conseils des trois communautés réunies pour le
fait des dits affranchissements et narré dans leur
délibération du dit jour douze octobre sujet de la
répartition à faire entre elles, il a été actuellement
arrêté qu’on s’en tiendra à ce qui a été convenu à ce
sujet, que chacune des dites trois paroisses payeront le
capital au dernier vingt du montant de la servitude
et des différents droits auxquels elle est en particulier
tenue, en conformité quant à la dite servitude des
cottes sommaires de chaque paroisse, et de la somme
portée par l’état que lesdits Seigneurs doivent en avoir remis
au bureau de la péréquation lors d’icelle et le surplus
du montant de la dites dépense sera réparti sur

\marginnote{\scriptsize Page 35}

les dites trois paroisses, et chacun suivant le bilan
de la taille royale, à la quelle chacune est taxée
mais, chacune des dites trois paroisses supportera en
son particulier, les frais de ses députés et autres
nécessaires à faire pour la réparation n’y ayant que
les sommes auxquelles arriveront le prix des acquisitions
des fiefs qui sont communs en elles, et s’étendent
sur leur territoire réciproque, ainsi que la finance
pour les droits réservés à la couronne qui doivent se
répartir entre elles, de la manière ci-dessus par la-
quelle répartition le député aura tout pouvoir,

\begin{center}
19
\end{center}

et pour procéder à tout ce que dessus, et faire tout ce qu’il
conviendra pour parvenir aux dites acquisitions, les susnommés
syndic, conseillers et particuliers, excédant les deux
tiers, à leur nom, et des autres communiers absents,
ont députés pour leur procureur spécial et général
une des qualités ne dérogeant à l’autre ni au contraire
Me Joseph fils du dit Joseph Bal notaire, et secrétaire
de la présente communauté d’Hauteluce \index[nom]{Bal@\textbf{Bal}  (Joseph, notaire et secrétaire de la communauté d’Hauteluce) }, habitant,
et originaire d’icelle, ici présent,, auquel les dits
comparants donnent  plein pouvoir d’agir, et passer
tous actes et contrats suivant et de la manière

\marginnote{\scriptsize Page 36}

ci-devant expliquée, pour les causes et cas y exprimés 
obliger avec clause de constitut les biens de la 
communauté et ceux des particuliers, ainsi que 
par le présent ils les obligent, et pour tous dépens, 
dommages, intérêts qui pourraient provenir de 
l’inexécution et généralement faire tout ce qui sera 
requis et nécessaire ici tenu pour exprimé, se pourvoir 
à tous tribunaux et office dont pourra ressortir la 
connaissance, approbation et vérification des contrats avec 
pouvoir de substituer procureur en son propre lieu 
et place, et qui aura le même pouvoir que lui et 
élisant domicile en la personne du substitut comme 
du constitué, avec promesse que font les dits constituants 
de relever leur dit procureur et ses substitués de toutes 
charges à supporter des dépens et frais qu’ils 
feront occasion du présent, et de s’en tenir à ces 
fins à la note qu’ils en donneront et pour ce qui est 
de leur vacation et transport dans les différents 
endroits, de même que pour leurs frais de voiture 
et nourriture le tout leur sera payé à la taxe 
qu’en faira  le seigneur intendant général aux mêmes

\marginnote{\scriptsize Page 37}

peines, obligations et constitution des biens de la 
communauté, et des particuliers d’icelle, les dits 
constituant par pacte exprès promettant de s’en 
tenir pour les dépenses à faire à l’occasion et en 
dépendance du contenu au présent à la note 
qu’en donnera le dit député de même qu’à celle 
qu’en donnera du temps de ses vacations, et de n’y 
contredire moyennant qu’elle soit assermentée 
de sa part sans autre forme et formalité à quoi 
ils renoncent par le présent par l’entière confiance 
qu’ils ont en sa probité. et clause requise fait 
et prononcé en présence de Maitre Jean-Baptiste Blanc 
notaire\index[nom]{Blanc@\textbf{Blanc} (Jean-Baptiste, notaire, témoin) }, et d’honorable Maxime Vibert Veched \index[nom]{Vibert Veched@\textbf{Vibert Veched} (Maxime, témoin) }tous deux 
de la paroisse de St Maxime témoins requis 
lesquels avec les dits. 




Joseph Bal-Fournier, \index[nom]{Bal-Fournier (Joseph) }
Louis Vionnet Fuapset, \index[nom]{Vionnet Fuapset (Louis) }
Révérent Baptandier vicaire, \index[nom]{Baptandier (Révérent, vicaire) }
Révérend. Reymond prêtre, \index[nom]{Reymond (Pierre, Rev. Sr. , prêtre et régent) }
Révérend. Girod Roux curé, \index[nom]{Girod Roux (François, Rev. Sr. , Prieur et curé d'Hauteluce) }
Louis Braisaz, \index[nom]{Braisaz (Louis) }
Louis Guiguet, \index[nom]{Guiguet (Louis) }
Jean-Pierre Burdillion, \index[nom]{Burdillion (Jean-Pierre) }
Louis Palluel-Blanc, \index[nom]{Palluel-Blanc (Louis) }
Claude Bal, \index[nom]{Bal (Claude) }
Joseph Cuvex Combaz, \index[nom]{Cuvex Combaz (Joseph) }
Louis Bron Bozon, \index[nom]{Bron Bozon (Louis) }
Jean-Baptiste Christiné, \index[nom]{Christiné (Jean-Baptiste) }
François Guiguet, \index[nom]{Guiguet (François) }
Pierre Bochet, \index[nom]{Bochet (Pierre) }

\marginnote{\scriptsize Page 38}

François Cuvex, \index[nom]{Cuvex (François) }
Jacques Malliard, \index[nom]{Malliard (Jacques) }
Jean-Baptiste Durand, \index[nom]{Durand (Jean-Baptiste) }
Jean François Bal pétré, \index[nom]{Bal Pétré (Jean François) }
Jean Guigue, \index[nom]{Guigue (Jean) }
Claude Durand, \index[nom]{Durand (Claude) }
François Martin, \index[nom]{Martin (François) }
François Pichol, \index[nom]{Pichol (François) }
Michel Baptandier Santiquet, \index[nom]{Baptandier Santiquet (Michel) }
Aymé Cuvex Combaz, \index[nom]{Cuvex Combaz (Aymé) }
Claude Malliard, \index[nom]{Malliard (Claude) }
Pierre Bal, \index[nom]{Bal (Pierre) }
Claude François Cabullin, \index[nom]{Cabullin (Claude François) }
Claude Braisaz, \index[nom]{Braisaz (Claude) }
Jacques Bal Pétré, \index[nom]{Bal Pétré (Jacques) }
Jean-Baptiste Vibert Bourgeois, \index[nom]{Vibert Bourgeois (Jean-Baptiste) }
Joseph Faure Couvillier, \index[nom]{Faure Couvillier (Joseph) }
Joseph Perroud, \index[nom]{Perroud (Joseph) }
Claude Duc Pithiod, \index[nom]{Duc Pithiod (Claude) }
Joseph Duret, \index[nom]{Duret (Joseph) }
Mammert Palluel Blanc, \index[nom]{Palluel Blanc (Mammert) }
Joseph Bal, \index[nom]{Bal (Joseph) }
Isidore Nicolas Gamin, \index[nom]{Gamin (Isidore Nicolas) }
André Ligeon, \index[nom]{Ligeon (André) }
Claude Favre Cotillet, \index[nom]{Favre Cotillet (Claude) }
Joseph Ligeon Ligeonnet, \index[nom]{Ligeon Ligeonnet (Joseph) }
Sylvestre Lachenal, \index[nom]{Lachenal (Sylvestre) }
Jean-Baptiste Ducretet, \index[nom]{Ducretet (Jean-Baptiste) }
Marie Joseph Guigoz, \index[nom]{Guigoz (Marie Joseph) }
François Carrey, \index[nom]{Carrey (François) }
Pierre Bonnet Ligeon, \index[nom]{Bonnet Ligeon (Pierre) }
Jean-Baptiste Ligeon Ligeonnet, \index[nom]{Ligeon Ligeonnet (Jean-Baptiste) }
Jean-Baptiste Ligeon, \index[nom]{Ligeon (Jean-Baptiste) }
André Favre Cotillet, \index[nom]{Favre Cotillet (André) }
Jacque DeDuc, \index[nom]{DeDuc (Jacque) }
Donat Bal, \index[nom]{Bal (Donat) }
Jean-Baptiste Braisaz,\index[nom]{Braisaz (Jean-Baptiste) }
et le dit Me Bal notaire\index[nom]{Bal@\textbf{Bal}  (Joseph, notaire et secrétaire de la communauté d’Hauteluce) } ont signé sur ma minute
et tous les autres constituant qui sont illettrés
de ce en quoi y ont fait leur marque, à moi notaire
sous-signé ait le présent expédié au dit Me Bal
après l’avoir faire insinuer au bureau du

\marginnote{\scriptsize Page 39}


Département de Beaufort [...]  trois cents
trente, payé trois livres pour le droit à former du
reçu du Sr.Blanc Secrétaire insinuateur
mon père\index[nom]{Blanc@\textbf{Blanc}  (Joseph, secrétaire insinuateur) }, du cinq […] mil sept cent soixante-
huit signé par me Blanc, notaire index[nom]{Blanc@\textbf{Blanc} (Michel, notaire et secrétaire de la communauté de St Maxime) }

(Signature manuscrite en bas de page)











\printindex[nom]
\newpage
\printglossary
\newpage
\printbibliography
\end{document}


